\section{Data split}
Neither the \gls{pdb} nor the \gls{crims} dataset can be split into training and validation sets using a simple random split. Two main characteristics of these datasets make such an approach inappropriate. The following describes the employed informed split strategy to avoid data leakage to enable a confident report in generalization performance.

First, the \gls{pdb} is highly redundant. The same protein is often deposited multiple times in the database, for example with different resolutions, slightly modified constructs, ligand bindings, or other experimental variations. At the time the data for this thesis was collected, the \gls{pdb} contained approximately 197,000 structure entries, but these corresponded to only about 89,000 unique proteins. Some proteins appear extremely frequently. For instance, the proteins R1AB\_SARS2 and LYSC\_CHICK had 2631 and 1258 associated structure files respectively. This is further depicted in \autoref{fig:most_common_proteins} in \autoref-{app:data}.


This redundancy has important implications for machine learning. As shown in \autoref{fig:proteinVsPEGComparison}, proteins tend to be highly selective with respect to their crystallization conditions. This selectivity arises from two main factors. First, the same protein often crystallizes reproducibly under similar conditions due to its structural and physicochemical properties. Second, once a successful crystallization condition has been reported in the \gls{pdb}, subsequent experiments frequently begin their condition search near previously successful conditions. This creates a strong bias linking specific proteins to specific crystallization conditions.

Although the protein name or identifier is not explicitly included as an input feature, the model can still indirectly learn protein identity through the protein sequence embedding. As illustrated in \autoref{fig:proteinEmbedding}, proteins form strong clusters in embedding space, indicating that the model could easily memorize protein-specific condition relationships. If a random split were used, it would therefore be very likely that the same protein—or even the same protein–condition pair—would appear in both the training and validation sets. Such data leakage would lead to overly optimistic performance estimates, as the model could simply recall previously seen patterns rather than learning generalizable relationships.

Avoiding leakage by splitting only by protein identity is insufficient. Although the determinants of crystallization conditions are not fully understood, properties such as structure, surface chemistry, and charge distribution strongly influence crystal formation \citep{Nanev2018}. These properties are implicitly encoded in the amino acid sequence, meaning similar sequences may lead to similar crystallization behavior. Thus, if related proteins appear in both training and validation sets, indirect leakage via sequence similarity can occur.

However, this assumption is debated. Studies such as \citet{Abrahams2019} and \citet{Liao2025} show that similar sequences can exhibit substantially different crystallization behavior, indicating that sequence similarity alone is not a reliable predictor. Nevertheless, to ensure rigorous validation and minimize potential leakage, sequence-based clustering is used when constructing data splits.

To address this issue, an additional level of separation based on sequence similarity clustering is required. A common strategy is to cluster proteins according to their sequence similarity and then ensure that proteins from the same cluster are assigned exclusively to either the training or the validation set. This ensures that the model cannot rely on closely related homologous proteins when making predictions.

For this purpose, sequence similarity clusters generated with MMseqs2 were used. Both the \gls{pdb} and \gls{crims} proteins were clustered at multiple sequence identity thresholds: 30, 40, 50, 70, 90, 95, and 100\% sequence identity. The clustering was performed using the easy-cluster algorithm of MMseqs2, which efficiently groups sequences into clusters based on pairwise sequence similarity. In this procedure, sequences are compared using fast k-mer based prefiltering followed by sequence alignment. Proteins are assigned to the same cluster if their sequence identity exceeds the specified threshold while satisfying coverage constraints (here using parameters c=0.8, cov-mode=0, and identity threshold ID=N). This ensures that sequences within a cluster share a minimum level of alignment coverage and sequence identity.

Lower identity thresholds produce broader clusters that group together more distantly related proteins, while higher thresholds result in tighter clusters containing only very similar or nearly identical sequences. Consequently, the 30\% identity clustering represents the most challenging evaluation scenario, as proteins in the validation set share only weak homology with proteins seen during training. In contrast, the 100\% identity clustering is the most permissive setting, where proteins in the validation set may be nearly identical to those present in the training set.

From an inference perspective, these splits correspond to different assumptions about the model's applicability. If the model performs well on the 30\% sequence identity split, it suggests that the model can infer crystallization conditions for proteins even when no closely related homologs are present in the training data. This would indicate strong generalization capability. On the other hand, strong performance on the 100\% identity split only demonstrates that the model can reproduce crystallization conditions for proteins that are essentially identical to proteins already present in the training set.

Evaluating the model across multiple sequence identity thresholds therefore provides a more comprehensive understanding of its generalization behavior, ranging from near-memorization scenarios to genuinely novel protein prediction tasks. \autoref{fig:train_test_split} depicts a schematic illustration. 
